
\section{Monotone ES-tree} \label{sec:EStrees}
Algorithm~\ref{alg:monotone_ES_tree} provides pseudocode for the monotone ES-tree algorithm, this is taken from~\cite{henzinger2014decremental}. 
\begin{algorithm}[H]
\caption{Monotone ES-tree}
\label{alg:monotone_ES_tree}


\SetKwFunction{initialize}{Initialize}
\SetKwFunction{delete}{Delete}
\SetKwFunction{increase}{Increase}
\SetKwFunction{insert}{Insert}
\SetKwFunction{updateLevels}{UpdateLevels}

\tcp{Internal data structures:}
\tcp{$ N (u) $: for every node $ u $ a heap $ N (u) $ whose intended use is to store for every neighbor $ v $ of $ u $ in the current graph the value of $ \lev (v) + w_{H} (u, v) $}
\tcp{$ Q $: global heap whose intended use is to store nodes whose levels might need to be updated}

\BlankLine

\procedure{\initialize{}}{
	Compute shortest paths tree from $ s $ in $ H $ up to depth $ L $\;
	\ForEach{$ u \in V $}{
		Set $ \lev (u) = \dist_{H} (s, u) $\;
		\lFor{every edge $ (u, v) $ in $ H $}{
			insert $ v $ into heap $ N(u) $ of $ u $ with key $ \lev(v) + w_{H} (u, v) $
		}
	}
}

\BlankLine

\procedure{\delete{u, v}}{
	\increase{$u$, $v$, $\infty$} 
}

\BlankLine

\procedure{\increase{u, v, $w (u, v)$}}{
	\tcp{Increase weight of edge $ (u, v) $ to $ w(u, v) $}
	Insert $ u $ and $v$ into heap $ Q $ with keys $ \lev(u) $ and $\lev(v)$ respectively\;\label{line:insert u}
	Update key of $ v $ in heap $ N(u) $ to $ \lev(v) + w(u, v) $ and key of $ u $ in heap $ N(v) $ to $ \lev(u) + w(u, v) $\;\label{line:update N after increase}
	$ \updateLevels{} $\;
}

\BlankLine

\procedure{\insert{$u$, $v$, $ w (u, v)$}}{
	\tcp{Increase edge $ (u, v) $ of weight $ w(u, v) $}
	Insert $ v $ into heap $ N(u) $ with key $ \lev (v) + w (u, v) $ and $u$ into heap $N(v)$ with key $ \lev (u) + w_{H}(u, v) $\; 
}

\BlankLine

\procedure{\updateLevels{}}{
	\While{heap $ Q $ is not empty}{
		Take node $ u $ with minimum key $ \lev (u) $ from heap $ Q $ and remove it from $ Q $\;
		$ \lev'(u) \gets \min_{v} (\lev (v) + w_{H} (u, v)) $\;
		\tcp{$\min_{v} (\lev (v) + w_{H} (u, v)) $ can be retrieved from the heap $ N(u) $. $\arg\min_{v} (\lev (v) + w_{H} (u, v)) $ is $u$'s parent in the ES-tree. }
	
		\If{$ \lev'(u) > \lev (u) $}{\label{line:check_for_level_increase}
			$\lev(u)\gets \lev'(u)$\;\label{line:level_increase}
			\lIf{$ \lev' (u) > L $}{
				$ \lev (u) \gets \infty $
			}

		\ForEach{neighbor $ v $ of $ u $}{
				update key of $ u $ in heap $ N(v) $ to $ \lev(u) + w_{H} (u, v) $\;
				insert $ v $ into heap $ Q $ with key $ \lev(v) $ if $Q$ does not already contain $ v $\;
			}
		}
	}
}
\end{algorithm}

\remove{
For the analysis of the monotone ES-tree we will use the following terminology. We say that an edge $(u,v)$ is \emph{stretched} if $\ell(u) > \ell(v) + w_H(u,v)$. We say that a node $u$ is \emph{stretched} if it is incident to an edge $(u,v)$ that is stretched. 


\begin{observation}[\cite{HKN2016}]\label{obs:simple_observations_monotone_ES}
The following holds for the monotone ES-tree:
\begin{enumerate}[label=(\arabic{*})]
\item \label{item: observation one} The level of a node never decreases.
\item \label{item: observation two} An edge can only become stretched when it is inserted.
\item \label{item: observation three} As long as a node is stretched, its level does not change. % Reason: can only change level in procedure reconnect
\item \label{item: observation four} For every tree edge $ (u, v) $ (where $ v $ is the parent of $ u $), $ \lev (u) \geq \lev (v) + w_{H} (u, v) $.
\end{enumerate}
\end{observation}
Observe that property~\ref{item: observation four} above implies that the returned distance estimate never underestimates the true distance.

\begin{lemma}[\cite{HKN2016}]\label{lm:ES-tree}
    For every $L \geq 1$, the total update time of a monotone ES-tree up to maximum level $L$ on a graph $H$ undergoing edge deletions, edge insertions, and edge weight increases is $O(\mathcal E(H)\cdot L+\mathcal W(H)+n)$, where $\mathcal E(H)$ is the total number of edges ever contained in $H$ and $\mathcal W(H)$ is the total number of updates to $H$.
\end{lemma}
}

\section{Maintaining Bunches and Clusters}\label{ap:bunches_and_clusters}
In this section we give a sketch of the algorithm of {\L}\k{a}cki and Nazari~\cite{LN2020} that allows us to maintain bunches and clusters needed for the decremental algorithm in \Cref{sec:TZ_2+eps}. 
One tool used is a decremental algorithm by Roditty and Zwick~\cite{RZ12} that maintains Thorup-Zwick clusters up to a bounded depth. This result can be summarized as follows:

\begin{lemma}[\cite{RZ12}]\label{lem:rz}
Let $0<p\leq 1$ and $d\geq 1$ be parameters. Given a dynamic graph $G=(V,E,w)$, that has edge deletions, weight increases, and edge insertions that do not decrease distances, we can maintain clusters (with sampling probability $p$) and bunches up to depth $d$ in total update time $\tilde O(|E(G)|d/p)$. Moreover, we maintain distances from each node to its cluster and its bunch, and these distances are monotone (non-decreasing).
\end{lemma}

This time dependence on $d$ is too expensive for many applications, and thus \cite{LN2020} improve the time for maintaining \textit{approximate} bunches and clusters by maintaining a hopset, combined with a well-known scaling tool. 

\paragraph{Hopsets.}
Hopsets, originally introduced by Cohen and Zwick~\cite{CohenZ01} in the context of parallel algorithms, are widely used in faster algorithms in decremental settings. We briefly define this concept and sketch its application. 
\begin{definition}
Let $G=(V,E,w)$ be a weighted undirected graph, and let $0<\epsilon <1$ be a parameter. A $(\beta,1+\epsilon)$-hopset is a graph $H=(V,E(H),w_H)$ such that for each $u,v\in V$ we have $d_G(u,v)\leq d^{(\beta)}_{G\cup H}(u,v) \leq (1+\epsilon) d_G(u,v)$, where $d^{(\beta)}_{G\cup H}(u,v)$ denotes the length of the shortest path between $u$ and $v$ in $G\cup H$ that uses at most $\beta$ hops. We say $\beta$ is the \emph{hopbound} of the hopset and $1+\epsilon$ is the \emph{stretch} of the hopset. 
\end{definition}


\begin{theorem}[\cite{LN2020}] \label{thm:hopset}
Let $0 < \rho < 1/2, 0< \epsilon <1$ be parameters. We can maintain a $(\beta,1+\epsilon)$-hopset $H$ for $G=(V,E,w)$ with hopbound $\beta =  O(\tfrac{\log(nW)}{\epsilon})^{2/\rho+1}$ in total update time $\tilde O (\tfrac{\beta\log^2(nW)}{\epsilon}mn^{\rho})$. Moreover, we can assume that the distances incurred by the additional edges of the hopset are monotone. Further we have that $H$ is of size $\tilde O(n^{1+\rho})$. 
\end{theorem}

By setting $\rho=\frac{\log \log n}{\sqrt{\log n}}$, we can maintain a hopset with hopbound $2^{\tilde{O}(\sqrt{\log n})}$ in $2^{\tilde{O}(\sqrt{\log n})}$ amortized time. We often choose this parameter setting when the graph is sparse for balancing the hopset maintenance time with the shortest path computation time.
By setting $\rho$ to a constant, we can maintain a hopset of polylogarithmic hopbound in time $O(mn^{\rho})$. We often use this parameter setting when we are maintaining distances from many sources and the graph is slightly denser (e.g.~$m=n^{1+\Omega(1)}$).

\paragraph{Scaling.}
The following scaling technique of Klein and Subramanian~\cite{KS97} changes a bounded hop distance problem into a bounded depth distance problem on a \textit{scaled graph}. Similar ideas have been widely used in dynamic algorithms, see e.g., \cite{Bernstein09,BR11,henzinger2014decremental,LN2020}. As we will see, the bounded depth works well together with the cluster and bunch maintaining algorithm of Roditty and Zwick~\cite{RZ12}
\begin{lemma}[\cite{KS97}]\label{lem:rounding}
Let $G=(V,E,w)$ be a weighted undirected graph. Let $R \geq 0$, $d \geq 1$, $\epsilon > 0$ be parameters.
We define a \emph{scaled graph} to be a graph $\textsc{Scale}(G, R, \epsilon, d)$ $:=  (V, E, \hat{w})$, such that $\hat{w}(e) = \lceil \tfrac{w(e)d}{R \epsilon} \rceil$.
Then for any path $\pi$ in $G$ such that $\pi$ has at most $d$ hops and weight $R \leq w(\pi) \leq 2R$ we have: 1) $\hat{w}(\pi) \leq \lceil 2d/\epsilon \rceil$, and 2) $w(\pi) \leq \tfrac{\epsilon R}{d} \cdot \hat{w}(\pi) \leq (1+\epsilon) w(\pi)$.
\end{lemma}

Hence given a hopset, we can `cover' the entire graph by scaled graphs, by considering separate scaled graphs $G_i$ for the distance scale $[2^i, 2^{i+1}]$ for $i=1,\dots, \log(nW)$. In order to combine this with a hopset algorithm, we need to scaled the graph after adding edges corresponding to a sequence of $2^i$-\textit{restricted hopsets}, i.e.~a graph that satisfies the hopset condition only for pairs within distance at most $2^i$. Slightly more formally, we define a scaled graph
as follows: $G^j := \textsc{Scale}(G \cup \bigcup_{r=0}^{j} H_{r}, 2^j, \epsilon, O(\beta)$, where $H_{r}$ is a $2^r$-restricted hopset. The final hopsets of \cite{LN2020} are the union of all these restricted hopsetes, however for maintaining distances efficiently we need to run the algorithm in Lemma \ref{lem:rz} on all these scaled graphs simultaneously and take the minimum estimate over all the scales. Importantly, to handle edge insertions they use monotone ES trees which implies that the distances obtained from this algorithm is monotone. This is important for us in bounding the running time of our decremental data structures.


%We can now use hopsets to maintain clusters in Theorem \ref{lem:rz} for all distances $d$ by removing the dependence $d$ in the running time (replaced with time dependence on $\beta$).
The following lemma shows that using the specific hopsets by \cite{LN2020}, we can maintain approximate single-source distances. 
%\begin{lemma}[\cite{LN2020}]\label{lm:hopset_scaledgraph}
%There is a decremental\footnote{Here by decremental we mean edge weights may only increase but there may be edge insertions.} $(1+\epsilon_1, \beta)$-hopset $H_i$ with properties described in theorem \ref{thm:hopset} that approximately maintains $\beta$-hop distance for all $u,v\in V$ with $d(u,v)\leq 2^{i+1}$. Moreover, for $i= 1 \leq i \leq \log (nW)$, $\textsc{Scale}(H_i,2^i,\epsilon,d)\cup G_i$ is also a monotone graph (i.e., distances are non-decreasing).
%\end{lemma}
%The fact that distances in $\textsc{Scale}(H_i,2^i,\epsilon,d)\cup G_i$ are non-decreasing over time allows us to maintain monotone ES-trees on it in time $O((m+E(H_i)))\beta)$, and hence $(1+\epsilon_1)$-SSSP. It also allows us to apply the bunch and cluster maintaining algorithm of Lemma~\ref{lem:rz}.


%To add: define a scaled graph $G_i$ for distance scale $[2^i, 2^{i+1}]$ in $G$.
%\begin{corollary}
%We can maintain clusters and bunches with parameter $p$ on $G_i$ in time $\tilde{O}(\frac{m\beta}{p})$.
%\end{corollary}
%\begin{proof}
%Combine scaling with the fact that we have hopset and plug in RZ theorem.
%\end{proof}

\lmbunch*
\begin{proof}[Proof sketch.]
First, using the algorithm of \cite{LN2020}, we can maintain a $(1+\epsilon, \beta)$-hopset $H$ of size $\tilde{O}(m+n^{1+\rho})$ in update time $\tilde{O}(mn^{\rho} \beta)$ with hopbound $\beta =\tilde{O}(\frac{1}{\rho\epsilon})^{2\rho}$, which is consisted of a sequence of $2^i$-restricted hopsets $H_1,...,H_{\log nW}$. We then use the framework of \cite{LN2020} to maintain \textit{approximate} clusters and hence approximate bunches on the scaled graphs. Since for all $u,v \in V$ in $G \cup H$ there is a path with $\beta$-hops that with length at most $(1+\epsilon)d_G(u,v)$, they show that it is enough to maintain $O(\beta/\epsilon)$-hop bounded bunches/clusters on scaled graphs $G_1,...,G_{\log n W}$. We consider $(4\beta/\epsilon)$-\textit{bounded bunches} $B_{G_i}(v)= \{ w: d_{G_i}(w,v) < d_{G_i}(v, p_i(v))\}$ where $p_i(v)$ is the closest sampled node in $G_i$ and $d_{G_i}(v, p_i(v)) \leq 4\beta/\epsilon$. The scaling ensures that for each node $v$ there exists $1 \leq i \leq \log nW$ such that $p_i(v)$ satisfies the property $d_{G_i}(v, p_i(v)) \leq 4\beta/\epsilon$. Let $d$ be the distances obtained by scaling back the estimates corresponding to $B_{G_i}(v) \cup p_i(v)$. There exists $i$ such that  $\tilde{p}(v) =p_i(v)$ satisfies $d_G(v,p(v)) \leq d(v,\tilde{p}(v)) \leq (1+\epsilon) d_G(v,p(v))$. Also by choosing the appropriate error parameters (to account for the scaling and hopset errors) in Theorem 14 of \cite{LN2020} we can ensure that for any $u \in \tilde{B}(v)$ we have $d_G(u,v) \leq d(u,v) \leq (1+\epsilon) d_G(u,v)$.
Consider any $u \in \hat{B}(v)$. We have $d(u,v) \leq (1+\epsilon)d_G(u,v) \leq d_G(v,p(v)) \leq d(v, \tilde{p}(v))$.
Hence we can see that the bunch $B_{G_i}(v) \subseteq \hat{B}(v)$ and thus $\tilde{B}(v):=\cup_{1 \leq i \leq \log nW} B_{G_i}$ contains all the nodes in $\hat{B}(v)$.

Finally since the algorithm of \cite{LN2020} (that is based on \cite{RZ12}) maintains bunches of size $O(\log n/p)$ on $O(\log nW)$ scaled graphs, we also have the bound on the size of the bunch.
\end{proof}



%\ttodo{add explicit APSP if we have time}
% \section{Explicit \texorpdfstring{$(2+\epsilon)$}{}-APSP}
% In this section, we show how to change our algorithm \DynAPSPtwo with constant query time in to an algorithm that explicitly maintains APSP in the same total update time. 
% \url{https://dl.acm.org/doi/pdf/10.1145/2344422.2344425}






